\documentclass[../../../include/open-logic-section]{subfiles}

\begin{document}

\olfileid{sth}{spine}{rank}
\olsection{秩}

既然我们已经把阶段定义为 $V_\alpha$，并且知道每个集合都是某个阶段的子集，我们就可以定义集合的\emph{秩}。直观上，集合 $A$ 的秩就是 $A$ 首次形成的时刻。更精确地说：

\begin{defn}\ollabel{defnsetrank}
对于每个集合 $A$，$\setrank{A}$ 是满足 $A \subseteq V_\alpha$ 的最小序数 $\alpha$。
\end{defn}
\begin{prop}\ollabel{ranksexist}
	对任意 $A$，$\setrank{A}$ 都存在。
\end{prop}
\begin{proof}
	留作练习。
\end{proof}
\begin{prob}
	证明 \olref[sth][spine][rank]{ranksexist}。	
\end{prob}
\noindent 
秩的良序性使我们能够证明一些重要的结果：
\begin{prop}\ollabel{valphalowerrank}
对任意序数 $\alpha$，$V_\alpha = \Setabs{x}{\setrank{x} \in \alpha}$。
\end{prop}

\begin{proof}
如果 $\setrank{x} \in \alpha$，那么 $x \subseteq V_{\setrank{x}} \in
V_\alpha$，所以 $x \in V_\alpha$，因为 $V_\alpha$ 是 potent 的（对子集封闭，这里多次用到了 \olref[Valphabasic]{Valphabasicprops}）。反之，如果 $x \in V_\alpha$，那么 $x \subseteq V_\alpha$，所以 $\setrank{x} \leq \alpha$；而通过简单的超限归纳可知 $x \notin V_\alpha$。 
\end{proof}
\begin{prob}
	完成 \olref[sth][spine][rank]{valphalowerrank} 中提到的简单超限归纳。
\end{prob}

\begin{prop}\ollabel{rankmemberslower}
如果 $B \in A$，那么 $\setrank{B} \in \setrank{A}$。
\end{prop}

\begin{proof}
由 \olref{valphalowerrank} 可得 $A \subseteq V_{\setrank{A}} = \Setabs{x}{\setrank{x} \in
\setrank{A}}$。
\end{proof}
\noindent
利用这一事实，我们可以建立如下结果，从而通过某种归纳形式来证明关于\emph{所有集合}的命题：

\begin{thm}[$\in$-归纳原理] 
对于任意公式 $\phi$：
\[
	\forall A((\forall x \in A)\phi(x) \lif \phi(A)) \lif \forall A \phi(A).
\]
\end{thm}

\begin{proof}
我们证明其逆否命题。假设 $\lnot \forall A \phi(A)$。根据超限归纳（\olref[ordinals][basic]{ordinductionschema}），存在一个秩最小的非 $\phi$ 对象；即存在一个 $A$ 使得 $\lnot \phi(A)$ 且 $\forall x(\setrank{x} \in \setrank{A}
\lif \phi(x))$。现在，如果 $x \in A$，则由 \olref{rankmemberslower} 知 $\setrank{x} \in \setrank{A}$，因此 $\phi(x)$；即 $(\forall x \in
A)\phi(x) \land \lnot \phi(A)$。
\end{proof}\noindent 对这个强有力的结果，有一种非形式的直观解释。称 $\phi$ 是\emph{遗传的}，如果当一个集合的所有元素都是 $\phi$ 时，该集合本身也是 $\phi$。那么 $\in$-归纳表明：如果 $\phi$ 是遗传的，那么每个集合都是 $\phi$。

为了（暂时）收尾关于秩的讨论，我们将证明几个前文已多次预告的命题。 

\begin{prop}\ollabel{ranksupstrict}
$\setrank{A} = \supstrict_{x \in A}\setrank{x}$。
\end{prop}

\begin{proof}
令 $\alpha = \supstrict_{x \in A}\setrank{x}$。由 \olref{rankmemberslower}，$\alpha \leq \setrank{A}$。但如果 $x \in A$，则 $\setrank{x} \in \alpha$，于是由 \olref{valphalowerrank} 知 $x \in V_\alpha$，因此 $A \subseteq V_\alpha$，即 $\setrank{A} \leq \alpha$。故 $\setrank{A} = \alpha$。
\end{proof}

\begin{cor}\ollabel{ordsetrankalpha}
对任意序数 $\alpha$，$\setrank{\alpha} = \alpha$。
\end{cor}

\begin{proof}
由超限归纳，假设对所有的 $\beta \in \alpha$ 都有 $\setrank{\beta} = \beta$。那么由 \olref{ranksupstrict} 知 $\setrank{\alpha} = \supstrict_{\beta \in
\alpha}\setrank{\beta} = \supstrict_{\beta \in \alpha}\beta = \alpha$。
%	First note that $\setrank{\alpha} \neq \beta$ for any $\beta \in
%	\alpha$. For otherwise we would have $\beta = \setrank{\beta} \in
%	\setrank{\alpha} = \beta$ by \olref{rankmemberslower}, i.e.,
%	$\beta \in \beta$, a contradiction. 
%
%	Now note that $\alpha \subseteq V_\alpha$. For $\alpha =
%	\Setabs{\beta}{\beta \in \alpha}$ by
%	\olref[sfr][ordinals][basic]{ordissetofsmallerord}. And if
%	$\beta \in \alpha$ then $\beta \subseteq V_{\setrank{\beta}} =
%	V_\beta \in V_\alpha$, hence $\beta \in V_\alpha$, by
%	\olref[sfr][spine][Valphabasic]{Valphabasicprops} twice. 
\end{proof}

最后，这里给出 \olref[foundation]{sec} 末尾所承诺结果的一个简短证明，即 $\ZFminus$ 能证明条件句 $\emph{正则性} \Rightarrow \emph{基础公理}$。（注意，由于“秩”的概念和 \olref{rankmemberslower} 可以在 $\ZFminus + \text{正则性}$ 中给出——正如本节开头所提——因此在本证明中可以使用它们。）

\begin{prop}[在 $\ZFminus + \text{正则性}$ 下工作]\ollabel{zfminusregularityfoundation} 基础公理成立。
\end{prop}

\begin{proof}
取定 $A \neq \emptyset$，并选取 $A$ 中一个秩最小的元素 $B$。如果 $c \in B$，则由 \olref{rankmemberslower} 知 $\setrank{c} \in \setrank{B}$，因此根据 $B$ 的选取，$c \notin A$。
\end{proof}

\end{document}